%==============================================================================
% appendices/C_dependent_proofs.tex
%==============================================================================
\section{Proofs for Dependent Extensions}
\label{appx:dependent-proofs}

This appendix proves the results stated in \cref{sec:dependent-cdmcs} and
\cref{sec:mrv-spectral}. Each proof environment is labelled by the result it
establishes.

\begin{proof}[Proof of \cref{thm:dep-cdmc-unbiased}]
Let
\[
  E_i(x)=\{S_N>x,\ R_i(X)=1\},
\]
where \(R_i\) is the selected-maximum rule of \cref{ass:tie}. The events
\(E_i(x)\) are disjoint and their union is \(\{S_N>x\}\). Conditional on
\(X_{-i}\), the event \(E_i(x)\) is equivalent to \(X_i\) exceeding both the
sum threshold and the selected-maximum threshold. In the continuous case this
is
\[
  X_i>(x-S_{-i})\vee M_{-i}=T_i(x).
\]
For discrete laws, replacing \(>M_{-i}\) by the tie-rule endpoint gives the
same identity with an adjusted threshold. Hence
\[
  \Pb(E_i(x))
  =\E\!\left[\Pb(X_i>T_i(x)\mid X_{-i})\right]
  =\E\!\left[p_i(T_i(x);X_{-i})\right].
\]
Summing over \(i\) proves the identity and the unbiasedness of
\(Z^{\mathrm{dep}}(x)\). No independence step is required.
\end{proof}

\begin{remark}[Population identity versus estimator]
\label{rem:pop-vs-est}
\Cref{thm:dep-cdmc-unbiased} contains two distinct statements. The
population-level identity
\(\mu(x)=\sum_i\E[p_i(T_i(x);X_{-i})]\) holds without any computability
assumption on the kernels. The estimator statement
\(\E Z^{\mathrm{dep}}(x)=\mu(x)\) holds provided the kernels are evaluated
exactly per replicate; if the kernels are evaluated by a noisy or approximate
routine, an additional bias term appears that is controlled by the kernel
approximation error.
\end{remark}

\begin{proof}[Proof of \cref{prop:dep-envelope-bre}]
By \cref{thm:dep-cdmc-unbiased}, \(\E Z^{\mathrm{dep}}(x)=\mu(x)\). The
selected-maximum threshold satisfies \(T_i(x)\ge x/N\) almost surely (the
argument is the same as in the independent CdMC proof in
\cref{appx:independent-proofs}). By \cref{ass:conditional-envelope}, for all
$t\ge x/N$,
\[
  p_i(t;X_{-i})\le K_{i,x}b_i(t),
\]
and therefore
\[
  Z^{\mathrm{dep}}(x)
  \le \sum_{i=1}^N K_{i,x}\, b_i(x/N).
\]
The second-moment envelope in \cref{ass:conditional-envelope} gives
\[
  \E\!\left[Z^{\mathrm{dep}}(x)^2\right]=O(\mu(x)^2),
\]
hence
\[
  \Var Z^{\mathrm{dep}}(x)\le \E Z^{\mathrm{dep}}(x)^2=O(\mu(x)^2),
\]
which is bounded relative error.
\end{proof}

\begin{proof}[Proof of \cref{thm:latent-shock-tail}]
Set \(Y_k=q_kZ_k\). Independence of the shocks transfers to the signed shock
contributions \(Y_k\). Their right-tail constants are
\[
  \Pb(Y_k>x)
  =\Pb(q_kZ_k>x)
  \sim |q_k|^\alpha\{p_k^+\1(q_k>0)+p_k^-\1(q_k<0)\}\Gb(x)
  =c_k(q_k)\Gb(x).
\]
Applying \cref{thm:sum-equivalence} to \(\sum_kY_k=q^\top Z\) yields the
first-order tail claim. Applying the independent CdMC construction and BRE
proof of \cref{appx:independent-proofs} to the \(K\)-dimensional shock vector
gives the estimator claim, with \(K\) replacing \(N\) in the BRE bound.
\end{proof}

\begin{proof}[Proof of \cref{thm:block-reduction}]
By hypothesis the blocks are independent, so the block sums \(Y_1,\ldots,Y_K\)
are independent. By assumption \(\Pb(Y_k>x)\sim d_k\Gb(x)\) and the block-level
single-big-jump conditions hold. \Cref{thm:sum-equivalence} applied to
\(\sum_kY_k\) gives
\[
  \Pb(S_N>x)=\Pb\!\left(\sum_kY_k>x\right)
  \sim \sum_k\Pb(Y_k>x)
  \sim \!\left(\sum_kd_k\right)\Gb(x).
\]
\end{proof}

\begin{proof}[Proof of \cref{thm:mrv-linear-risk}]
By multivariate regular variation,
\[
  \frac{\Pb(X/x\in A)}{\Hb(x)}\to \nu(A)
\]
for every Borel set \(A\) bounded away from the origin with
\(\nu(\partial A)=0\). Take \(A=A_\ell=\{z:\ell(z)>1\}\); this set is bounded
away from the origin and admissible by the boundary-continuity hypothesis of
\cref{ass:mrv-loss}. Then
\[
  \Pb\{\ell(X)>x\}=\Pb(X/x\in A_\ell)\sim \Hb(x)\nu(A_\ell).
\]
For the polar representation,
\[
  \nu(A_\ell)=\int\!\!\int
  \1\{r\ell(\theta)>1\}\,\alpha r^{-\alpha-1}\dd r\,S(\dd\theta).
\]
If \(\ell(\theta)\le 0\) the inner integral is zero. If \(\ell(\theta)>0\) the
inner integral equals \(\int_{1/\ell(\theta)}^\infty \alpha r^{-\alpha-1}\dd r
=\ell(\theta)^\alpha\). Integrating over \(\theta\) gives the stated spectral
integral.
\end{proof}

\begin{proof}[Proof of \cref{prop:radial-cdmc}]
Condition on \(\Theta\). If \(\ell(\Theta)\le 0\), then \(\ell(R\Theta)>x\) is
impossible for \(x>0\). If \(\ell(\Theta)>0\), the event is
\(R>x/\ell(\Theta)\). Thus
\[
  \Pb\{\ell(X)>x\mid\Theta\}
  =\Fb_R\!\left(x/\ell(\Theta)\mid\Theta\right)\,\1\{\ell(\Theta)>0\}.
\]
Taking expectations yields the population identity and proves unbiasedness of
\(Z^{\mathrm{spec}}(x)\).
\end{proof}

\begin{proof}[Proof of \cref{prop:spectral-bre}]
Let \(Y=\ell(\Theta)_+\). The envelope assumption gives
\[
  Z^{\mathrm{spec}}(x)\le K\Hb(x)Y^\alpha,
\]
hence
\[
  \E\!\left[Z^{\mathrm{spec}}(x)^2\right]
  \le K^2\Hb(x)^2\E Y^{2\alpha}.
\]
The MRV linear-risk asymptotic gives \(\mu(x)\sim \Hb(x)\E Y^\alpha\) under the
same normalization. Therefore
\[
  \limsup_{x\to\infty}\frac{\Var Z^{\mathrm{spec}}(x)}{\mu(x)^2}
  \le K^2\frac{\E Y^{2\alpha}}{(\E Y^\alpha)^2}<\infty,
\]
which is bounded relative error provided \(\E Y^{2\alpha}<\infty\).
\end{proof}
